%% kapa: exam=1 round=32 year=2021 period=1 form=A subject=civil_law subject_ko=민법 seq=11 number=11
%% source: Q-Net 감정평가사 기출문제 (cst00309) artlSeq=5208397 file=extracted/1008/5208397/A형/2021년도 제32회 감정평가사 1차 1교시 A형.pdf
%% conversion: kordoc 4.24.51 -> tools/gen1.py (auto); review=proofread-sonnet; flags=none
\begin{question}{11}
통정허위표시에 의하여 외형상 형성된 법률관계를 기초로 하여 ‘새로운 법률상 이해관계를 맺은 제3자’에 해당하지 않는 자는? (다툼이 있으면 판례에 따름)
\begin{choices}
  \item 가장전세권에 관하여 저당권을 취득한 자
  \item 가장소비대차에 기한 대여금채권을 양수한 자
  \item 가장저당권 설정행위에 기한 저당권의 실행에 의하여 목적부동산을 경락받은 자
  \item 가장의 채권양도 후 채무가 변제되지 않고 있는 동안 채권양도가 허위임이 밝혀진 경우에 있어서의 채무자
  \item 가장소비대차의 대주(貸主)가 파산한 경우, 파산자와는 독립한 지위에서 파산채권자 전체의 공동의 이익을 위하여 직무를 행하게 된 파산관재인
\end{choices}
\end{question}
\answer{4}
